Agent memory retrieval requires a selector that delivers relevant items and an
accounting of the cost to build them. We study temporal dominance --- a failure mode
in the deployed hybrid selector where a fixed query-independent freshness weight
displaces lexically relevant older records in favour of fresher but irrelevant ones.
We introduce CTLS (Contrastive Tier-Based Lexical Selection), a novel selector that
partitions candidates by lexical overlap tier before applying temporal scoring,
provably eliminating temporal dominance by construction. The partition satisfies a
Non-Dominance Property: for any two candidates where one has positive Jaccard overlap
with the query and the other has zero overlap, CTLS always ranks the positive-overlap
item first, regardless of freshness gap, access frequency, or importance values, and
regardless of the weight assigned to the temporal term within the relevant-tier score.
On a stratified LongMemEval-S subset of 84 memory episodes under three token budgets,
BM25 outperforms the deployed hybrid by 36.1, 29.4, and 33.7 percentage points
respectively (all Holm-corrected; $p \le 0.0022$ throughout). The hybrid\_no\_time
ablation recovers most of this gap, confirming the temporal term as the primary cause.
CTLS formalizes this fix while preserving temporal decay as a within-tier tiebreaker
rather than discarding it globally. The mechanism inequality derived in this paper
shows exactly when temporal dominance occurs, explains why budget growth does not help
the hybrid, and establishes hybrid\_no\_time as a conservative lower bound on CTLS.
CTLS requires no model calls, no embedding model, and no learned parameters; it is a
zero-cost selector substitution that leaves construction cost and write amortization
unchanged. We report the formal Non-Dominance Property with proof sketch, controlled
evidence across three token budgets with cluster bootstrap inference, and lifecycle
cost under memory reuse.
